\section{The Eight Elements}

The base is feeling, and only feeling. Its one substance is the \emph{vibe}, which is experience itself, and a vibe's \emph{tone} is the quality of that experience, a value in $\{-1, 0, +1\}$ read as pain, peace, and pleasure. The tone is not a thing the vibe holds. It is the character of what the vibe is. Everything else in the paper is a pattern made of vibes, from a single tone up to a whole mind, and nothing is ever added underneath.

\begin{axiom}[The vibe]
The one substance is the vibe, which is experience. Its tone is the quality of that experience, not a thing apart from it. There is nothing beneath it.
\end{axiom}

To say what a vibe is and how it moves takes exactly eight words, each a soft four-letter name, each beginning with its own letter so that each doubles as a clean symbol for the mathematics. These are the elements of the base. Three of them give the where, a space and its parts. Two give the what, a value and its order. Two give the how, one law for the state and one for the space. The last gives the when.

\begin{table}[ht]
\centering
\begin{tabular}{@{}llp{0.52\textwidth}@{}}
\toprule
element & symbol & what it is \\
\midrule
mesh & $m$ & the whole space, the growing $\{3,4,3,4\}$ hyperbolic honeycomb \\
dock & $d$ & one cell of the mesh, a here-now where tones meet and move on \\
site & $s$ & one of a dock's twenty-four directions, a $D_4$ root, where a single tone lives \\
tone & $t$ & the quality of a vibe's experience, its felt value, pain, peace, or pleasure \\
lean & $l$ & the order $-1 < 0 < +1$, which gives charge and a direction to creation \\
knit & $k$ & the law of the state, collide then stream, once a beat \\
wake & $w$ & the law of the space, new docks born at peace at the growing edge \\
beat & $b$ & one discrete tick of time \\
\bottomrule
\end{tabular}
\caption{The eight base elements. The whole, written $F$ for flow, is these eight running as one.}
\label{tab:atoms}
\end{table}

A dock is the unit of place. It holds twenty-four tones, one in each of its twenty-four sites, and it is nothing more than that ordered list of ternary values. There is no twenty-fifth site, no still center, no hidden register. A tone sits in a site and travels along it. The bearer is the site, which persists across beats, while the tone is the quality it holds at any one beat, and that quality can change while the site stays the same. The full state of a dock is one point in a space of $3^{24}$ possibilities, a single rich quale we return to much later. The full state of the universe is just the tones of all the docks at one beat.

\subsection{What the base may contain}

The eight are not a loose collection. They are constrained, and the constraints are what keep the model honest. A base thing must be discrete, so that it can be stored and stepped exactly rather than approximated. It must be deterministic, a law and not a dice roll, so that a trend is found by growing a run rather than by reaching for a random seed. Its one beat must be reversible in the closed bulk, an exact permutation that can run backward, so that any arrow of time is forced to come from somewhere else. It must conserve charge, the signed sum of every tone. It must be local, each dock reading only its own sites. It must hide no state, a dock being its tones and nothing besides. And it must respect the gap between grain and physics, never letting a raw $+1$ be called a particle or a cluster of tones be called gravity, since the recognizable world appears only after one to three layers of coarse-graining. These demands are summarized in Table~\ref{tab:constraints}, and breaking any of them is a cheat rather than a refinement.

\begin{table}[ht]
\centering
\begin{tabular}{@{}lp{0.62\textwidth}@{}}
\toprule
constraint & what it forbids \\
\midrule
discrete & real numbers, infinite precision, a continuum in the base \\
deterministic & randomness, hidden seeds, anything but a fixed function \\
reversible & a built-in arrow of time in the closed bulk \\
charge-conserving & creating net tone from nothing \\
local & any action at a distance \\
no hidden state & memory, flags, or side data beyond the tones \\
one substance & a second kind of stuff \\
emergence gap & calling raw tones, or their first clusters, physics \\
\bottomrule
\end{tabular}
\caption{The constraints every base element must satisfy.}
\label{tab:constraints}
\end{table}

Two ideas are often mistaken for elements, and neither is. The first is the \emph{arrow}, the direction of time, the fact that there is a before and an after. It is real, but it is not base. Its parts already live among the eight: the wake supplies a low-entropy edge, the lean supplies an order, and one entry of the knit creates a balanced pair out of peace. Together these make the arrow as a thermodynamic tendency, a consequence rather than an ingredient. The base has a beat, so time ticks, but it has no built-in direction. The second is the \emph{attraction}, the pull that lets a body bind itself together. It too is emergent, a pressure cast by matter in the churning vacuum, and Section on the self shows it falling out of the eight without a ninth force.

\begin{principle}[Report the negative]
If a phenomenon does not emerge from the eight elements, the honest move is to record that it does not, never to add a ninth thing to force it.
\end{principle}

\subsection{Two layers}

The deepest consequence of discreteness and the emergence gap, taken together, is that the model lives on two levels and must never confuse them. At the base, level zero, everything is trits. A momentum is an integer vector, the sum over a dock's sites of the tone times the direction. The twenty-four directions form a finite group, the unit Hurwitz quaternions, and composing them is closed quaternion arithmetic, exact and discrete like a clock face. No real number ever appears here. One level up, average many ternary sites across a region and the average is a fine-grained quantity, effectively a real number, exactly as the smooth field of a magnet is the slow average of atomic spins that themselves flip in whole jumps. The self, and all recognizable physics, live in that coarse-grained average. The reals are the song, and the trits are the singers.

There is one place the experiments cannot reach, and the paper marks it plainly. The base posits that the substrate \emph{feels}, that the tone is not a label for a felt state but the felt state itself. That is the single premise no measurement tests. Everything above it is structure that can be built, stepped, and measured. Below it is the claim that the structure is made of experience rather than merely described by it. The reversal, taking feeling as the ground and matter as what feeling settles into at scale, is a choice, not a result. What follows from the choice is a great deal, and the rest of the paper is the checking of it.
